\documentclass[11pt]{article}
\usepackage[T1]{fontenc}
\usepackage{lmodern}
\usepackage[margin=1.04in]{geometry}
\usepackage{microtype}
\usepackage{amsmath,amssymb,amsthm,mathtools,booktabs}
\usepackage[hidelinks]{hyperref}
\hypersetup{pdftitle={Sharp Hardy coefficient norms in arbitrary degree}}
\newtheorem{theorem}{Theorem}
\newtheorem{proposition}[theorem]{Proposition}
\newtheorem{lemma}[theorem]{Lemma}
\newtheorem{corollary}[theorem]{Corollary}
\newcommand{\D}{\mathbb D}
\newcommand{\T}{\mathbb T}
\newcommand{\Dif}{\mathcal D}
\newcommand{\coef}[1]{[z^{#1}]}
\newcommand{\Ent}{\operatorname{Ent}}
\allowdisplaybreaks[1]
\title{Sharp Hardy coefficient norms in arbitrary degree}
\author{}
\date{27 November 2026}
\begin{document}
\maketitle
\begin{abstract}
For every \(n\geq1\) and \(0<p<1\), we express the sharp
Taylor coefficient norm \(C(n,p)\) as a maximum over a finite
list recovered from a single explicit characteristic determinant.
For fixed \(n\), the norm has only finitely many possible
transitions between real analytic branches over the full
exponent range. We obtain an explicit small-exponent uniqueness
range of order \(n^{-1}\). At \(p=1\) we give an exact arithmetic
description of the limiting logarithmic slope, prove that every
limiting extremal factor has simple unit-circle zeros, and
compute the outward velocity of every zero along an extremizing
sequence. A second-variation criterion proves coefficient positivity
in the range \(p\le2/n\), and the limiting entropy operator has
spectral gap exactly two.
\end{abstract}

\section{The coefficient problem}

Let \(m\) denote normalized arc measure on \(\T\), and put
\[
 \|f\|_p=\sup_{0<r<1}\left(\int_\T|f(rz)|^p\,dm(z)\right)^{1/p},
 \qquad
 C(n,p)=\sup_{\substack{f\in H^p(\D)\\\|f\|_p\leq1}}|[z^n]f|.
\]
Brevig and Saksman \cite{BrevigSaksman} formulated the sharp
coefficient problem and gave exact results in low degree.
Brevig and S{\o}vik \cite{BrevigSovik} established outerness
of the optimizing functions. Earlier work obtained explicit
branches in degrees three through five and a general finite
cell-based representation \cite{Low,High,Fifth}. The present
paper uses the reflected stationarity and coefficient
elimination methods to obtain a determinant and quantitative
results uniform in the degree, building on the coefficient methods
of the Verifiable Mathematical Reasoning Corpus \cite{VMRC}. The associated endpoint problem
is the maximization of entropy over unit-root polynomial
densities.

\section{Boundary separation and stationary jets}

\begin{proposition}\label{prop:boundary}
Fix \(n\ge1\), \(0<p<1\), and \(s=2/p\). The norm \(C(n,p)\)
is attained. Every extremizer of quasi-norm one has the form
\[
 f=\eta Q^s,\qquad
 \deg Q=n,\quad \|Q\|_2=1,\quad Q(0)>0,\quad
 Q(z)\ne0\ \text{for }|z|\le1,\quad |\eta|=1 .
\]
The power uses the analytic logarithm of \(Q\) which is real at zero.
\end{proposition}

\begin{proof}
The Hardy pointwise estimates give a normal family on every smaller
disk. A locally uniform limit of a maximizing sequence still has
quasi-norm at most one, by integration on each smaller circle.
The coefficient converges and the maximum is positive, since \(z^n\)
is admissible. Scaling shows every extremizer has quasi-norm one.
Brevig--S{\o}vik's outer theorem \cite[Thm.~1]{BrevigSovik} applies
to the bounded linear functional \(\coef n\). For convenience, its
concavity argument is as follows. Write \(f=IO\), where \(I\) is
inner and \(O\) is outer. If \(I\) is nonconstant, set
\(u=2\Re I-c\) on \(\T\), with
\(c=\int_{\T}2\Re I\,|f|^p\,dm\). Then \(u\) is nonconstant and has
mean zero under \(|f|^p\,dm\). The analytic function
\(g=O(1+I^2-cI)\) belongs to \(H^p\), and \(g=uf\) on \(\T\).
For small real \(t>0\), strict concavity of \(x^p\) gives
\(\|f\pm tg\|_p^p<1\), since \(1\pm tu>0\).
Averaging the two coefficient values contradicts extremality.
Thus \(f\) is outer. Here \(\Re I\) cannot be constant on \(\T\)
unless \(I\) is constant, and the outer boundary modulus is positive
almost everywhere.

By rotations assume \(a_0=f(0)>0\) and \(a_n=\coef n f=C(n,p)>0\).
Vary \(f e^{t g}/\|f e^{t g}\|_p\), first with an analytic polynomial
\(g\) and then with \(ig\). Differentiating the norm is justified
by boundedness of \(g\) on \(\T\). Stationarity gives
\[
 \coef n(fg)=a_n\int_{\T}|f|^p g\,dm .
\]
With \(g=z^k\) for every \(k\ge0\), Fourier uniqueness gives
\[
 h(z)=a_n+\sum_{k=1}^n
       \bigl(a_{n-k}z^{-k}+\bar a_{n-k}z^k\bigr)
     =a_n |f(z)|^p \quad\text{on }\T\text{ a.e.}
\]
The Fej\'er--Riesz theorem factors \(h=a_n |Q|^2\) with
\(Q\) free of zeros in \(\D\), \(Q(0)>0\), \(\|Q\|_2=1\), and
degree exactly \(n\) because \(a_0>0\). The normalized outer function
with boundary modulus \(|Q|^{2/p}\) is \(Q^s\), so \(f=Q^s\).

If \(U\) has degree at most \(d\), write \(U^{*d}(z)=
z^d\overline{U(1/\bar z)}\).
Let \(R=z^n h=a_n QQ^{*n}\), let \(\Dif=z\,d/dz\), and set
\(E=R\Dif\log f=s a_n Q^{*n}\Dif Q\).
The polynomial \(R\) is self-reciprocal of degree \(2n\), and
\(R=f\bmod z^{n+1}\), meaning their Taylor coefficients agree through
order \(n\). Consequently the first \(n+1\) coefficients of \(E\)
are \(k a_k\), \(0\le k\le n\). The product rule and
\((\Dif Q)^{*n}=nQ^{*n}-\Dif(Q^{*n})\) give
\[
 E-E^{*2n}=s(\Dif R-nR).
\]
Set \(\beta=1-1/p=1-s/2\). The first \(n+1\) coefficients of
\((E+E^{*2n})/2\) are \((n-\beta(n-k))a_k\).
Reflecting the remaining coefficients shows that, on \(\T\),
\begin{equation}\label{eq:boundary-identity}
 z^{-n}\frac{E+E^{*2n}}{2}=nh-\beta D_h,
\end{equation}
where \(D_h\) is the outward radial derivative of the harmonic
extension of \(h\), the Fourier multiplier \(|k|\).

Suppose \(Q(\zeta)=0\) for some \(\zeta\in\T\). Then \(E(\zeta)=0\)
by its factored form and \(E^{*2n}(\zeta)=0\) by reflection.
After rotation to \(\zeta=1\), write \(Q(z)=(z-1)B(z)\).
Applying the multiplier to
\(h=a_n(2-z-z^{-1})|B|^2\) at 1 gives
\(D_h(1)=-2a_n\|B\|_2^2<0\), since
\(2|k|-|k+1|-|k-1|\) vanishes except at \(k=0\), where it
equals \(-2\). Equation \eqref{eq:boundary-identity} would now give
\(0=-\beta D_h(\zeta)\), which is impossible. This proves strict
boundary separation.
\end{proof}

There are two useful forms of stationarity. Put \(a=s-1>1\),
\(q_0=Q(0)\), and \(P=Q/q_0\). With \(a_n>0\) as in the proof,
the leading coefficient \(d=\coef n P\) is positive because the
constant term of \(R=a_n QQ^{*n}\) is \(q_0^s>0\).
The congruence \(R=f\bmod z^{n+1}\) gives
\begin{equation}\label{eq:jet}
 P^a\equiv P^{*n}/d\pmod{z^{n+1}},\qquad
 a_n=\frac{q_0^{a-1}}{d},\qquad
 q_0=(1+|[z]P|^2+\cdots+|\coef nP|^2)^{-1/2}.
\end{equation}
When rotating the argument without the \(a_n>0\) normalization,
\(P^a\equiv P^{*n}/\bar d\bmod z^{n+1}\) and
\(\coef n Q^s=q_0^{a-1}/\bar d\). These identities also follow
from \eqref{eq:jet} by rotation.
The boundary argument also applies to a formal solution of the
jet with \(P(0)=1\) and \(d\ne0\). To see this, multiply the
jet by \(P\), take the formal power \(f=P^{a+1}\) at zero, and
use \(R=PP^{*n}/\bar d\) after choosing \(d>0\) by rotation.
The proof of \eqref{eq:boundary-identity} then uses only Taylor
coefficients and the finite trigonometric polynomial
\(z^{-n}R=|P|^2/d\) on \(\T\). It therefore still excludes a
unit-circle root. This observation controls limits of solutions
when the exponent varies.

\section{A determinant formula for every order}\label{sec:all-n-det}

Fix \(n\geq1\) and \(0<p<1\), and put \(a=2/p-1\). Write
\[
 P(z)=1+r_1z+\cdots+r_{n-1}z^{n-1}+dz^n,
 \qquad r_n=d\in\mathbb R,\qquad r_0=1.
\]
Write \(k_0=\lfloor(n-1)/2\rfloor\) and retain only the real and
imaginary parts of \(r_1,\ldots,r_{n-k_0-1}\), together with \(d\).
Thus the number of real coordinates \(X=(x_1,\ldots,x_m)\) is
\(m=n\) for odd \(n\) and \(m=n+1\) for even \(n\).
For \(1\leq k\leq k_0\), reconstruct
\(r_{n-k}=d\,\overline{A_k(X)}\), where
\[
 A_k(X)=[z^k]P^a
   =\sum_{j=1}^{k}\binom aj [z^k](P-1)^j .
\]
The reconstruction uses only the retained coordinates because
\(A_k\) uses \(r_1,\ldots,r_k\). The positive sign for \(d\) will
be imposed at the end. Define the residual polynomial
\[
 J(X)=\sum_{k=k_0+1}^n
           \left|dA_k(X)-\overline{r_{n-k}}\right|^2 .
\]
Assign weight \(j\) to both real coordinates of the retained
coefficient \(r_j\) and weight one to \(d\). Denote the weights
by \(w_i\). The reconstructed \(r_{n-k}\) has weighted degree
at most \(k+1\leq n-k\), so each \(r_j\) has weighted degree
at most \(j\). Taylor weight in \(A_k\) is \(k\), hence
\(\deg_w A_k\leq k\) and \(\deg_w J\leq 2n+2\).
Set
\[
 b_i=2\lfloor(n+1)/w_i\rfloor+1,\qquad
 M=\prod_{i=1}^m b_i
   =(2n+3)\prod_{j=1}^{\lfloor n/2\rfloor}
        \bigl(2\lfloor(n+1)/j\rfloor+1\bigr)^2 .
\]
For a formal parameter \(h\),
use the following \(m\) reduction rules:
\begin{equation}\label{eq:matrix-reduction}
 X_i^{b_i}=-\frac{1}{(b_i+1)h}\frac{\partial J}{\partial X_i}
       \quad (1\leq i\leq m).
\end{equation}
The ordinary monomial basis consists of the \(M\) monomials
\(X^\nu\) with \(0\leq\nu_i<b_i\). With \(H\) an independent
indeterminate, let \(B(H,h)\) be the multiplication matrix of
\(\sum_{i=1}^m H^{i-1}X_i\) in this basis, using
\eqref{eq:matrix-reduction}. Expand
\[
 \chi(T,H,h)=\det(TI_M-B(H,h))
\]
as a polynomial in \(h^{-1}\). Let \(E(T,H)\) denote the
coefficient of its largest power of \(h^{-1}\).

\begin{theorem}\label{thm:matrix-all-n}
The polynomial \(E\) is nonzero and has \(T\)-degree at most
\(M\). Take its monic factors
\[
 T-\sum_{i=1}^m x_i H^{i-1}\qquad (x_i\in\mathbb R)
\]
in \(\mathbb R(H)[T]\). From each such factor keep the vector
\(X=(x_1,\ldots,x_m)\) if \(J(X)=0\), \(d>0\), and the
associated \(P\) is zero-free on \(\overline{\mathbb D}\).
This gives a nonempty finite list with at most \(M\) elements,
and
\[
 C(n,p)=
 \max_{P\ \mathrm{on\ the\ list}}
      \frac{1}{d\|P\|_2^{a-1}} .
\]
All coordinates of the list are algebraic over \(\mathbb Q(a)\).
\end{theorem}

\begin{proof}
Every normalized extremizer has an associated stationary factor
\(P\) whose last coefficient can be made positive by an argument
rotation. Its coefficient vector lies in \(J=0\).
Conversely any vector passing the three tests gives an admissible
\((P/\|P\|_2)^{2/p}\) with coefficient modulus
\(1/(d\|P\|_2^{a-1})\), by multiplying the reflected jet by \(P\). The linear projection to retained coordinates and the polynomial reconstruction are inverse homeomorphisms on the jet set.

The formal jet equations exclude \(d=0\) by the equation for
\(k=n\). The reflected boundary identity, applied to the formal
jet at zero, excludes a unit-circle root even without assuming
\(P\) is stable. Thus both the sign of \(d\) and the number of
zeros of \(P\) in the disk are constant on each connected
component of the real algebraic set \(J=0\). A component with
positive \(d\) and all zeros outside the closed disk is compact:
the reciprocal-root representation bounds
\(\lvert r_k\rvert\leq\binom nk\), and a connected component
of the closed set \(J=0\) is closed. The real algebraic set has
finitely many connected components.

Let \(K\) be any such compact component. Choose a bounded open
neighborhood \(V\) of \(K\) whose closure meets \(J=0\) only in
\(K\). This is possible since the other finitely many components
are closed. The minimum of \(J\) on \(\partial V\) is strictly
positive. For all sufficiently small positive \(h\),
\[
 J_h(X)=J(X)+h\sum_i X_i^{b_i+1}
\]
then achieves its minimum on \(\overline V\) at an interior
point \(X(h)\). Indeed the perturbation is nonnegative, and the
value at a fixed point of \(K\) tends to zero. These points
satisfy \eqref{eq:matrix-reduction}. A subsequence as \(h\downarrow0\)
converges to some \(x\in K\), since \(J(X(h))\to0\).

Over the field with \(h\) inverted, the left-hand monomials in
\eqref{eq:matrix-reduction} are the leading monomials for any
order refining weighted degree, since the right-hand sides
have weighted degree at most \(2n+2-w_i<w_i b_i\). They are pairwise coprime, so the
polynomials form a Gr\"obner basis by the product criterion.
This proves the quotient basis and the asserted polynomial
dependence on \(h^{-1}\) and \(H\). For every specialization
\(h>0\), evaluation at any solution of the critical equations is
an eigenvector of the transpose multiplication matrix, with
eigenvalue \(\sum_i H^{i-1}X_i(h)\). Multiply the characteristic
equation by the largest power of \(h\) occurring in the
definition of \(E\) and pass to the indicated subsequence.
For every real \(H\) we get
\[
 E\left(\sum_i H^{i-1}x_i,H\right)=0.
\]
It follows as a polynomial identity in \(H\) that the factor
associated to \(x\) occurs. Distinct vectors yield distinct
linear factors, so there are at most \(M\) vectors on the list.

The candidate value is constant along any continuously
differentiable stationary path. For a direct verification, write
\(s=a+1\) and \(F=U^{-s/2}[z^n]P^s\) with \(U=\|P\|_2^2\).
The jet gives \(F=U^{1-s/2}/d\), which is real. Differentiating
the first expression for \(F\) and using
\[
 \Re\bigl([z^n]P^a\dot P\bigr)
   =\frac1d\Re\bigl([z^n]P^*\dot P\bigr)=\frac{\dot U}{2d}
\]
shows that \(\dot F=0\).
Semialgebraic connected components are joined by piecewise
continuously differentiable paths. Thus the vector chosen by the
determinant in a maximizing component has the maximizing value.

Finally, the factor identity can be evaluated at \(m\) distinct
rational values of \(H\) for which \(E(T,H)\) is nonzero as a
polynomial in \(T\). Only finitely many values are excluded.
Each evaluated root is algebraic over \(\mathbb Q(a)\); interpolation
gives the same property for its \(m\) coefficients.
\end{proof}

For a fixed algebraic \(p\), factorization, real-root isolation,
and a Schur test are finite exact operations on algebraic
numbers. One way to recover the real linear factors is to use
the \(m\) rational evaluation points in the proof, take all
tuples of real roots of the specialized polynomials, interpolate,
and test the polynomial identity in \(H\). At rational \(p\)
the candidate powers can be compared algebraically. At irrational
algebraic \(p\), comparing two candidates with different
\(\|P\|_2^2\) terminates by the Gelfond--Schneider theorem and
certified logarithm comparisons. Equal norms reduce immediately
to comparing the leading coefficients.

\paragraph{Size of the determinant.}
Writing \(k=\lfloor n/2\rfloor\), every retained weight satisfies
\((n+1)/j\geq2\), so
\(2\lfloor(n+1)/j\rfloor+1\leq 5(n+1)/(2j)\).
Thus \(M\leq(2n+3)(5(n+1)/2)^{2k}/(k!)^2\).
For \(k\geq1\), use \(k!\geq(k/e)^k\) and
\(((n+1)/(2k))^{2k}\leq e^2\) to get
\[
 M\leq e^2(2n+3)(5e)^n .
\]
The same inequality holds at \(n=1\). In particular the
univariate degree bound is singly exponential in the order.

\paragraph{Degree bounds.}
If \(X\) is a sample, then
\[
 [\mathbb Q(a,x_1,\ldots,x_m):\mathbb Q(a)]\leq M.
\]
Indeed, each embedding of the sample field into an algebraic
closure over \(\mathbb Q(a)\) sends the linear form
\(\sum_i x_iH^{i-1}\) to a distinct root of \(E(T,H)\).
When \(p=r/t\) is reduced and rational,
\(C(n,p)^r=d^{-r}\|P\|_2^{-2(t-r)}\) belongs to the field of
a maximizing sample. Consequently,
\[
 [\mathbb Q(C(n,p)):\mathbb Q]
       \leq rM.
\]

\paragraph{The first coefficient.}
For \(n=1\) we have \(J=(ad^2-1)^2\) and the reduction rule is
\(d^5=-(2a/3h)d(ad^2-1)\).
The characteristic polynomial is
\[
 T^5+\frac{2a}{3h}T(aT^2-1).
\]
Its leading pole coefficient selects the positive stable root
\(d=a^{-1/2}\), giving
\(C(1,p)=\sqrt a(1+1/a)^{-(a-1)/2}\).

\paragraph{A single univariate reconstruction polynomial.}
For each fixed \(a\), the two-variable determinant can be converted
to one polynomial in \(T\) without needing to factor it over
\(\mathbb R(H)\). Let \(S(T,H)\) be the primitive part in \(T\) of
its squarefree part over \(\mathbb Q(a)(H)\), with denominators
cleared. Terms depending on \(H\) alone are discarded. Choose a
rational \(H_0\) so that the leading coefficient and
\(\operatorname{Res}_T(S,\partial_T S)\) are nonzero at \(H_0\).
There exists such an integer \(H_0\) in
\(0\leq H_0\leq 2(m-1)M^2\). Indeed
\(\deg_H S\leq\deg_H E\leq (m-1)M\), since the multiplication
matrix entries have \(H\)-degree at most \(m-1\), and the product
of the leading coefficient and resultant has \(H\)-degree at
most \(2M\deg_H S\).

Isolate the real roots \(t\) of \(S(T,H_0)\), all simple. For each
root take the Taylor polynomial \(R_t(H)\) of order \(m-1\)
about \(H_0\) of the local root \(T=\tau(H)\) with \(\tau(H_0)=t\).
Its coefficients are rational functions of \(a,H_0,t\), found
by implicit differentiation of \(S(\tau(H),H)=0\). Keep only the
polynomials satisfying the identity \(E(R_t(H),H)\equiv0\).
Their coefficients in powers of \(H\) are precisely the desired
linear-factor vectors. Thus the determinant formula requires
one univariate root list of degree at most \(M\), followed
by rational reconstruction and admissibility tests.

\section{A spectral test for stationary factors}

Throughout this section put \(s=2/p>2\), \(a=s-1\), and
\(Q(z)=\sum_{j=0}^n q_jz^j\) with \(\|Q\|_2=1\).
When \(q_0\ne0\), the formal coefficient
\[
 F_s(Q)=[z^n]Q^s
\]
is defined using a local logarithm at zero. Its modulus is
independent of that choice. The formal powers and derivatives
below only require \(q_0\ne0\).

\begin{lemma}[Finite Schur interpolation]\label{lem:general-schur-interp}
For a jet \(b(z)=\sum_{j=0}^n b_jz^j\bmod z^{n+1}\), let
\(T_n(b)\) be the lower triangular Toeplitz matrix representing
multiplication by \(b\) on polynomials modulo \(z^{n+1}\).
The jet has a holomorphic extension to the disk with
modulus at most one if and only if \(T_n(b)\) is a contraction.
\end{lemma}

\begin{proof}
Necessity follows by compression of multiplication on \(H^2\).
For sufficiency, if \(|b_0|=1\), the first column norm forces
\(b_1=\cdots=b_n=0\), so a constant extension works. The case
\(n=0\) also has a constant extension. If \(|b_0|<1\), put
\[
 X=(T_n(b)-b_0 I)(I-\bar b_0 T_n(b))^{-1}.
\]
It is a contraction because
\[
 I-X^*X=(1-|b_0|^2)
 (I-b_0 T_n(b)^*)^{-1}(I-T_n(b)^*T_n(b))
 (I-\bar b_0 T_n(b))^{-1}.
\]
The first row and last column of \(X\) are zero. Its
submatrix on the last \(n\) rows and first \(n\) columns
is \(T_{n-1}(c)\), where the jet \(c\) is given by
\(zc=(b-b_0)/(1-\bar b_0 b)\bmod z^{n+1}\).
This matrix is a contraction. By induction \(c\) has a Schur
extension, and \((b_0+zc)/(1+\bar b_0 zc)\) extends \(b\).
\end{proof}

\begin{theorem}\label{thm:general-hessian-schur}
Every local maximum of the formal coefficient modulus
\(|F_s(Q)|\) on the \(H^2\) unit sphere at which \(q_0F_s(Q)\ne0\)
has degree \(n\) and is zero-free on the closed disk. More generally,
suppose \(q_0F_s(Q)\ne0\) and \(Q\) satisfies the first-order equations
\[
 Q^{s-1}\equiv F_s(Q)\,Q^{*n}\pmod {z^{n+1}}.
\]
Define the complex symmetric matrix \(H\) by
\[
 H_{ij}=[z^{n-i-j}](Q^{*n}/Q),\qquad 0\le i,j\le n,
\]
with negative coefficients interpreted as zero. If
\begin{equation}\label{eq:general-hankel-hessian}
 |u^T H u|\le \|u\|_2^2/a\qquad \text{whenever }q^*u=0,
\end{equation}
then \(Q\) is zero-free on the closed disk.
\end{theorem}

\begin{proof}
Write \(F=F_s(Q)\). At a nonzero local maximum, Lagrange
multipliers for \(|F_s|^2\) and homogeneity give
\(\partial_{q_j}F_s=sF\bar q_j\). The left side is
\(s[z^{n-j}]Q^{s-1}\), which gives the stated first-order
equations. In particular \(q_0^{s-1}=F\bar q_n\ne0\), so
\(Q\) has degree \(n\). For \(u\) with \(q^*u=0\), let
\(U(z)=\sum_j u_j z^j\) and take the real curve
\[
 Q_t=\frac{Q+tU}{(1+t^2\|U\|_2^2)^{1/2}}.
\]
The first-order equations make the first derivative of
\(F_s(Q_t)\) at zero equal to zero. They also give
\(
 Q^{s-2}\equiv F Q^{*n}/Q\bmod z^{n+1}.
\)
It follows that the second derivative of
\(\log|F_s(Q_t)|\) at zero is
\(s(a\Re(u^T H u)-\|u\|_2^2)\). It is nonpositive at a
local maximum. Replacing \(u\) by a unimodular multiple proves
\eqref{eq:general-hankel-hessian}.

We now only assume the first-order equations and the spectral
test. Multiplication by \(Q\) in the definition of \(H\) shows
\(Hq=\bar q\). Symmetry of \(H\) implies it maps \(q^\perp\)
into \(\bar q^\perp\).
For unit \(u,v\in q^\perp\), polarization of the symmetric
bilinear form and \eqref{eq:general-hankel-hessian} give
\[
 |u^T H v|
 \le \frac{\|u+v\|_2^2+\|u-v\|_2^2}{4a}=\frac1a .
\]
Consequently the operator norm of \(H\) on \(q^\perp\) is at
most \(1/a<1\), and \(H\) has norm one with a one-dimensional
norm-attaining subspace. Row reversal of \(H\) gives
\(T_n(Q^{*n}/Q)\). Lemma~\ref{lem:general-schur-interp} extends
the jet of \(Q^{*n}/Q\) to a Schur function \(b\).
The first \(n+1\) coefficients of \(bQ\) equal
those of \(Q^{*n}\), and
\[
 \|bQ\|_2\le\|Q\|_2=\|Q^{*n}\|_2 .
\]
Parseval therefore gives \(bQ=Q^{*n}\) on the disk.

If \(Q\) and \(Q^{*n}\) have a common factor of degree \(g>0\),
write \(b=A/B\) in reduced form with numerators and denominator
of degree \(n-g\). On the unit circle \(|b|=1\) except at most
finitely many points, since \(|Q^{*n}|=|Q|\). For each polynomial
\(V\) of degree at most \(g\), \(bBV=AV\) is a polynomial of
degree at most \(n\), and
\(\|T_n(b)BV\|_2=\|BV\|_2\).
This would give a norm-attaining subspace of dimension at least
\(g+1\), a contradiction. Thus \(Q\) and \(Q^{*n}\) are coprime.
The equation \(bQ=Q^{*n}\) excludes a root of \(Q\) in the open
disk, and a root on the circle would be common to the two
polynomials. This proves the assertion.
\end{proof}

The spectral condition is equivalently the matrix inequality
\[
 H^*H\ \preceq\ a^{-2}I+(1-a^{-2})qq^* .
\]
In fact \(Hq=\bar q\), and the polarization argument gives the
norm bound on \(q^\perp\) in one direction. The operator bound
gives the quadratic-form bound in the other.
Thus stationary points passing this algebraic test furnish
admissible values for the original coefficient problem, and
every extremizer passes the test.

\begin{theorem}\label{thm:general-small-positive}
For all \(n\ge1\) and \(s>n\) with \(s>2\), the sharp Hardy
coefficient is the maximum of \(|F_s|\) over the entire
\(H^2\) unit sphere, where \(|F_s|\) is extended by zero at
\(q_0=0\). The same conclusion, with the ordinary polynomial
power, holds for every integer \(s\ge3\) and every \(n\).
In either range every extremal factor can be made strictly
positive coefficientwise by rotation of the argument and
multiplication by a unimodular scalar. In particular this
conclusion holds for \(0<p<1\) and \(p\le 2/n\).
\end{theorem}

\begin{proof}
If \(s>n\), the finite binomial expansion has terms of the
form \(q_0^{s-k}\) times a polynomial in \(q_1,\ldots,q_n\),
where \(k\le n\). Each term tends uniformly to zero on the
sphere as \(|q_0|\to0\). The extended modulus is continuous.
A trial polynomial supported in degrees \(0,n\) gives a
positive value, so a global maximum has \(q_0F_s(Q)\ne0\).
Theorem~\ref{thm:general-hessian-schur} makes this maximizer
admissible for the original Hardy problem. The structural
Hardy reduction puts every Hardy maximizer on the same sphere,
so the two maxima coincide.

If \(s\ge3\) is an integer, every ordinary power \(Q^s\) with
\(\|Q\|_2=1\) is admissible. The two maxima are again equal by
the Hardy reduction. A maximizing \(Q\) must have \(q_0\ne0\)
by outerness of the corresponding Hardy extremizer.

In both ranges, the expansion has nonnegative real
monomial weights after making \(q_0>0\), and taking all
coefficient moduli gives a maximizing vector. Write \(S\)
for the support of this nonnegative maximizer, and \(J\)
for the support of \(Q^{s-1}\) through degree \(n\).
The first-order equations give \(J=n-S\). Positivity of the
linear term gives \(S\subseteq J\), and consequently
\(S=J=n-S\). The quadratic term has a positive coefficient
when its index lies through degree \(n\), so \(S\) is closed
under sums at most \(n\). Reflection makes it closed
under nonnegative differences. Since \(0,n\in S\), this
implies \(S=\{0,h,\ldots,n\}\) for a divisor \(h\) of \(n\).

For \(h>1\), the original power is a function of \(z^h\),
so its coefficient is at most \(C(n/h,p)\). This is strictly
smaller than \(C(n,p)\). Indeed multiplication of a lower
order maximizer by the appropriate power of \(z\) preserves
its Hardy norm and places its sharp coefficient in degree
\(n\); equality would contradict outerness.
It follows that \(h=1\) and all the moduli are positive.
In the original complex maximizer all monomial phases
must agree in the triangle inequality. With \(q_0>0\),
comparison of the monomials
\(q_jq_{n-j}q_0^{s-2}\) and
\(q_1q_{j-1}q_{n-j}q_0^{s-3}\) for \(2\le j\le n\)
gives \(\arg(q_j)=j\arg(q_1)\bmod 2\pi\).
The variable rotation then makes all coefficients positive.
\end{proof}

In the positivity range of Theorem~\ref{thm:general-small-positive},
one can use only the real retained coefficients in
Theorem~\ref{thm:matrix-all-n}. The same proof applies on that real
stationary locus, which contains a maximum after rotation.
The matrix order becomes
\[
 (2n+3)\prod_{j=1}^{\lfloor n/2\rfloor}
        \bigl(2\lfloor(n+1)/j\rfloor+1\bigr).
\]

\section{Finitely many analytic branches}

Use \(s=2/p>2\) as the real parameter. For a normalized factor
\(Q=P/\|P\|_2\), write \(U=\|P\|_2^2\) and
\(v=|[z^n]P|^2\). A stable stationary factor has value
\(v^{-1/2}U^{-(s-2)/2}\). A branch is called algebraic here
when every coordinate is a real analytic algebraic function
of \(s\) on its stated interval.

\begin{theorem}\label{thm:branches}
For each \(n\geq1\), there is a finite set of parameters in
\((2,\infty)\) such that, on every complementary open interval,
one algebraic branch \(P(s)\) attains \(C(n,2/s)\). On that
interval
\[
 \log C(n,2/s)=-\tfrac12\log v(s)-\tfrac{s-2}{2}\log U(s).
\]
The second derivative of \(\log C(n,2/s)\) there is algebraic.
\end{theorem}

\begin{proof}
The stable stationary jet locus in \((s,P)\) is semialgebraic
over \(\mathbb Q\). The jet involves only finite generalized
binomial polynomials. Stability is also semialgebraic, for
instance by excluding a root with real and imaginary coordinates
satisfying \(u^2+v^2\leq1\) by real quantifier elimination.
Take a finite cylindrical algebraic decomposition with \(s\)
as the first coordinate. Over each open interval of the
parameter partition, it gives finitely many semialgebraic
sections, one in each nonempty fiber cell. For example,
in each lifted interval choose a midpoint, or an endpoint
plus or minus one for a one-sided unbounded interval.
A further finite partition makes every section real analytic
and algebraic. Each fiber cell is path-connected by piecewise
smooth paths \cite{BCR}, and the stationary coefficient value
is constant along such paths \((s\text{ fixed})\).
Indeed, differentiating the coefficient along normalized \(Q\)
on the complex unit sphere and using the reflected jets gives
\[
 \frac{d}{dt}\log|[z^n]Q^s|
 =s\Re\sum_{j=0}^n\overline{q_j}\,\dot q_j=0.
\]
Here the gradient formula follows from
\([z^{n-j}]Q^{s-1}=([z^n]Q^s)\bar q_j\).
It follows that the finitely many sections attain all possible
stationary values, and an extremizer belongs to one of the cells.

For any pair of sections, twice their logarithmic difference
is of the form \(A(s)+(s-2)B(s)\), where \(A,B\) are logarithms
of positive algebraic functions. Its second derivative is
algebraic, so it is either identically zero or vanishes at
only finitely many points on the interval. If it is identically
zero, the difference is affine. Otherwise Rolle's theorem
bounds the number of zeros of a nonidentically-zero difference
by two more than the number of zeros of its second derivative.
Finitely many pairwise comparisons therefore give a finite
partition of the entire interval \(2<s<\infty\), on which a
single branch attains the maximum. The same differentiation
observation proves the claim about the second derivative.
\end{proof}

For a density \(W\) on \(\T\), write
\(\Ent(W)=\int_{\T} W\log W\,dm\), with \(0\log0=0\).
Define \(\tau=1/p-1=(s-2)/2\), and put
\(J_n(\tau)=\log C(n,1/(1+\tau))\).

\begin{proposition}\label{prop:slopes}
At every \(\tau>0\), the left and right derivatives of \(J_n\)
exist and are, respectively, the minimum and maximum of
\(\Ent(|Q|^2)\) over its normalized extremal factors.
At an algebraic \(\tau>0\), both derivatives are transcendental.
\end{proposition}

\begin{proof}
The feasible set of normalized powers can be taken to be all
\(Q=P/\|P\|_2\) with \(P(0)=1\), degree at most \(n\), and
reciprocal roots in the closed disk, padding the product with
zero reciprocal roots as needed. This is compact, as the
continuous image of \(n\) closed disks. Every such \(Q^s\)
is admissible, including when \(Q\) vanishes on the boundary,
and the outer and jet reduction shows that their maximum is
\(C(n,p)\). The Taylor coefficient of \(Q^s\) is analytic in
\(s\) at any fixed factor, and its derivative is continuous
jointly in \(s,P\) by the finite Taylor expansion at zero.
We can use the squared modulus to avoid any singularity
at a zero of the coefficient. The one-sided envelope rule
now follows by comparing with a fixed maximizing factor
for the lower bound in a right difference quotient, and
taking a convergent subsequence of nearby maximizing factors
for the upper bound. The left-sided inequalities are reversed.

At an extremizer, the reflected identity implies, for every
function \(g\) analytic on a neighborhood of the closed disk,
\[
 [z^n](Q^s g)=([z^n]Q^s)\int_\T |Q|^2g\,dm .
\]
This follows by multiplying the reflected jet by \(P\) and
comparing the first \(n+1\) coefficients. Take \(g=\log Q\),
using the analytic branch real at zero, which exists by
boundary separation. The derivative with respect to \(\tau\)
of \(\log|[z^n]Q^{2+2\tau}|\) at fixed \(Q\) is then
\(2\Re\int |Q|^2\log Q\,dm=\Ent(|Q|^2)\).

If \(\tau\) is algebraic, take a semialgebraic decomposition
of the stable stationary fiber over the real algebraic numbers.
Its finite cells have algebraic samples, the value is constant
on each cell, and hence the union of maximizing cells is
semialgebraic over that field. The entropy of their factors is
\[
 \Ent(|Q|^2)=A(P)-\log U,\qquad
 A(P)=\frac{2}{U}\Re\int_\T |P|^2\log P\,dm.
\]
The real part in the integral uses only the first \(n\) Taylor
coefficients of \(\log P\), so \(A(P)\) is a rational function
of the real and imaginary coefficients of \(P\) over
\(\mathbb Q\). The extrema of the entropy on the compact set
of maximizing factors each occur at an algebraic sample
of the following semialgebraic critical locus. In a finite
smooth semialgebraic cell decomposition of that set, require
the differential of \(A-\log U\) to vanish on the tangent
space to the cell. This is a semialgebraic condition because
its differential is \(dA-dU/U\). The same path-connected
cell argument makes entropy constant on each connected
component of that critical locus inside the cell. Choosing
an algebraic sample on the appropriate component gives
the claimed algebraic data. Finally \(U>1\) because the
degree-\(n\) coefficient is nonzero. By the
Hermite--Lindemann theorem, \(\log U\) is transcendental
for algebraic \(U>1\), so the resulting two entropy values
are transcendental.
\end{proof}

\section{An explicit uniqueness range in every order}

\begin{theorem}\label{thm:explicit-small-p}
Let \(n\geq2\) and \(0<p\leq(2048n^2)^{-1}\).
Then the \(n\)th Hardy coefficient has exactly one unit-norm
extremizer modulo constant phases and argument rotations.
Write \(s=2/p\), \(a=s-1\), and \(\rho=\sqrt{n/s}\).
Its factor \(P\), normalized by \(P(0)=1\) and \([z]P>0\),
is the unique stable positive solution of the reflected jet
in the region
\[
 P(z)=1+\rho\sum_{j=1}^n x_j z^j,\qquad
 3/4\leq x_1,\quad 0\leq x_j\leq1\ (j\geq2),\quad
 \sum_{j=1}^n x_j^2\leq2.
\]
With \(d=[z^n]P\), its sharp value is
\(\|P\|_2^{-(a-1)}/d\).
The factor and value depend real analytically on \(p\) in the
interior of the displayed range.
\end{theorem}

\begin{proof}
Let \(Q(z)=\sum_{j=0}^n q_jz^j\) be the unit-\(H^2\) factor
of any extremizer, with \(q_0>0\).
The dilation \(Q(e^t z)/\|Q(e^t z)\|_2\) is admissible
for real \(t\) near zero by strict boundary separation.
The derivative of the log modulus of its \(n\)th coefficient
after raising it to \(s\) gives
\begin{equation}\label{eq:uniform-euler}
 \sum_{j=1}^n j |q_j|^2=n/s=\rho^2.
\end{equation}
Consequently
\(\sum_{j=1}^n |q_j|\leq\rho\sqrt{H_n}<\sqrt{1-\rho^2}
\leq q_0\), where \(H_n=\sum_{j=1}^n 1/j\).
Thus the coefficientwise modulus of \(Q\) gives a stable factor.
Since \(\binom{s}{k}>0\) for \(1\leq k\leq n\), this
replacement cannot decrease the coefficient modulus.
It suffices first to study the maximum on nonnegative factors.
For such a maximizing factor put \(x_j=q_j/(\rho q_0)\).
By \eqref{eq:uniform-euler} we have \(x_j\leq1\) for \(j\geq2\)
and \(\|x\|_2^2\leq1/(1-\rho^2)<2\).

We prove \(x_1>c=3/4\) by comparison with the trial factor
having \(x_1=1\) and all other coordinates zero. The ratio
of the normalization factor of \(P\) to that of the trial,
after taking the power \(s\), is at most
\[
 \left(\frac{1+\rho^2}{1+\rho^2 x_1^2}\right)^{s/2}
 \leq \exp\bigl(n(1-x_1^2)/2\bigr)
 \qquad(0\leq x_1\leq c).
\]
In a monomial of \([z^n](\sum x_jz^j)^{n-\ell}\),
the number \(m_1\) of factors \(x_1\) is at least \(n-2\ell\).
Write \(h=c\exp((1-c^2)/2)\). Since
\(x\exp((1-x^2)/2)\) is increasing on \([0,1]\), the
monomial times the displayed normalization ratio is at most
\[
 h^{m_1}\exp((n-m_1)/2)\leq h^n(e/h^2)^\ell .
\]
The number of compositions of length \(n-\ell\) is at most
\(n^\ell/\ell!\), and
\[
 \frac{\binom{s}{n-\ell}\rho^{n-\ell}}{\binom{s}{n}\rho^n}
 \leq\left(\frac{\rho}{1-\rho^2}\right)^\ell
 \leq(4\rho/3)^\ell .
\]
Summing proves that the objective ratio to the trial is bounded
by \(h^n\exp(4en\rho/(3h^2))<1\). Indeed
\(-\log h>1/16\), \(4e/(3h^2)=64e^{9/16}/27<128/27\),
and \(\rho\leq1/(64\sqrt n)\) with \(n\geq2\).
For the first elementary inequality one may use the
four-term alternating lower bound for \(\log(1+1/3)\);
the second uses \(\log2\geq 2/3\).
This rules out \(x_1\leq c\).

Consider the convex set \(\Omega\) given in the theorem.
Every polynomial on it is stable, since
\(\rho\sum x_j\leq\rho\sqrt{2n}<1\).
The logarithm of the normalized coefficient there is, up
to a constant,
\begin{equation}\label{eq:uniform-log-objective}
 n\log x_1-\frac s2\log(1+\rho^2\|x\|_2^2)+\log R(x),
 \qquad
 R(x)=\frac{[z^n]P^s}{\binom{s}{n}\rho^n x_1^n}\geq1.
\end{equation}
We bound the upward curvature of the correction \(\log R\).
For \(j,k\geq2\), let
\(M_j=\tfrac43 n(2\rho)^{j-1}\). Then
\begin{equation}\label{eq:uniform-replacement}
 0\leq R_j\leq M_jR,\qquad 0\leq R_{jk}\leq M_jM_kR .
\end{equation}
To prove the first inequality, use the expansion on a
multiplicity vector \(m\) with \(\sum i m_i=n\).
If \(k_0=\sum m_i\), its unnormalized coefficient is
\((s)_{k_0}\rho^{k_0}/\prod m_i!\), with \((s)_k\) the
falling factorial. After differentiating one \(x_j\), compare
the resulting term to the undifferentiated term obtained by
replacing this \(x_j\) by \(j\) copies of \(x_1\).
The ratio is bounded by
\[
 \frac{n^j}{x_1^j((s-n)\rho)^{j-1}}
 \leq \tfrac43 n(2\rho)^{j-1}.
\]
The replacement is injective on multiplicity vectors.
Two replacements give the second estimate, including when
the differentiated coordinates coincide. The same comparison
applies at zero coordinates by continuity.

Each monomial of \(R\) has the form
\(A\prod_{j\geq2}(x_j/x_1^j)^{m_j}\). Thus
\(x_1R_1=-\sum_{j\geq2}jx_jR_j\).
Put \(T_i=\sum_{j\geq2}j^iM_j\),
\(S_0=T_0\), \(S_1=\sum_{j\geq2} jx_jM_j\),
and \(S_2=\sum_{j\geq2} j(j+1)x_jM_j\).
Differentiating the identity and using
\eqref{eq:uniform-replacement} gives
\[
 \frac{|R_{11}|}{R}\leq c^{-2}(S_2+S_1^2),\qquad
 \frac{|R_{1j}|}{R}\leq c^{-1}M_j(j+S_1)\quad (j\geq2).
\]
In particular the sum of absolute entries of \(D^2R/R\)
is at most \(c^{-2}(S_2+S_1^2)+2c^{-1}(T_1+S_1S_0)+S_0^2\).
Geometric sums with \(2\rho\leq1/32\) give
\[
 S_0\leq3n\rho,\qquad T_1\leq6n\rho,\qquad
 S_1\leq6n\rho,\qquad S_2\leq18n\rho .
\]
Using \(c^{-2}<2\) and \(2c^{-1}<3\), we obtain
\[
 D^2\log R
 \preceq (54n\rho+135n^2\rho^2)I\prec 0.64n I .
\]
The Hessian of the normalization term in
\eqref{eq:uniform-log-objective} is bounded above by
\(-n(1-v)/(1+v)^2 I\) with \(v=\rho^2\|x\|_2^2\),
and therefore by \(-0.99nI\) on \(\Omega\).
Since \(n\log x_1\) is concave, the entire log objective is
strictly concave. Its maximum on \(\Omega\) is unique.

At that maximum, the first coefficient is positive and all
coefficients of \(P^a\) through degree \(n\) are positive:
the term using only \(x_1\) is positive in every order.
The reflected jet forces all coefficients of \(P\) to be
positive. It also determines the objective and says that its
gradient is zero. Conversely, any stable positive reflected
jet in the region is a stationary point of
\eqref{eq:uniform-log-objective}. For this, note that the
derivative of \([z^n]P^s\) in direction \(u\) is
\(s[z^n](P^a u)\); the jet turns this into the multiple of
the norm derivative that cancels the normalization.
Strict concavity gives uniqueness of that point.
The negative definite Hessian and the analytic implicit
function theorem give analytic dependence on \(p\).

Finally, for an arbitrary complex extremizer, its moduli
give this same positive maximizing factor. In the equality
case of the coefficientwise triangle inequality, compare
the monomial \(q_1^n q_0^{s-n}\) with the monomial in
\(q_jq_1^{n-j}\), with its appropriate power of \(q_0>0\).
Every monomial has a positive real weight before the phases
are introduced. Their equality of phases gives
\(\arg q_j=j\arg q_1 \pmod{2\pi}\).
An argument rotation thus yields precisely the unique
positive factor.
\end{proof}

\begin{corollary}\label{cor:uniform-iteration}
Under the assumptions of Theorem~\ref{thm:explicit-small-p},
let \(\Phi\) denote the log objective in
\eqref{eq:uniform-log-objective}, and let
\(\Pi_\Omega\) be Euclidean projection onto its convex region.
Starting at \(x^{(0)}=(1,0,\ldots,0)\), the iteration
\[
 x^{(m+1)}=\Pi_\Omega\left(x^{(m)}
       +\frac{1}{4n}\nabla\Phi(x^{(m)})\right)
\]
converges to the maximizing factor coordinates \(x^\star\).
For every \(m\geq0\),
\(\|x^{(m)}-x^\star\|_2\leq2\sqrt2(11/12)^m\).
\end{corollary}

\begin{proof}
The same derivative bounds give
\[\|\nabla R\|_2/R\leq c^{-1}S_1+S_0\leq11n\rho.\]
Consequently
\(\|D^2\log R\|_{\mathrm{op}}<0.67n\), after adding
the bound for its gradient rank-one term to the absolute
entry-sum bound for \(D^2R/R\).
The normalizing Hessian has operator norm at most \(n\),
and that of \(n\log x_1\) is at most \(16n/9\).
Together with the upper bounds in the theorem, the eigenvalues
of \(D^2\Phi\) on \(\Omega\) therefore lie in
\([-4n,-n/3]\).
Integrating the Hessian along a segment shows that
\(x\mapsto x+\nabla\Phi(x)/(4n)\) is Lipschitz with constant
at most \(11/12\). Euclidean projection is nonexpansive,
and the stationary maximizer is a fixed point.
The bound follows by iteration and \(\operatorname{diam}
\Omega\leq2\sqrt2\).
\end{proof}

\section{A uniform norm comparison}

\begin{theorem}\label{thm:uniform-comparison}
Let \(s=2/p\geq4096n\), and put \(\rho=\sqrt{n/s}\).
Define the degree-one trial value
\[
 M(n,s)=\binom{s}{n}\rho^n(1-\rho^2)^{(s-n)/2}.
\]
Then
\[
 M(n,s)\leq C(n,p)\leq M(n,s)\exp\!\left(\frac{27n^2}{4s}\right),
 \qquad
 e^{-n^2/s}\leq
 \frac{n!\,C(n,p)}{(ns/e)^{n/2}}
 \leq e^{8n^2/s}.
\]
The estimates hold for every \(n\geq1\).
\end{theorem}

\begin{proof}
The unit-\(H^2\) polynomial
\(\sqrt{1-\rho^2}+\rho z\) is stable and gives \(M(n,s)\)
after raising it to \(s\).
For the upper bound take an extremal unit-\(H^2\) factor \(Q\)
with \(q_0>0\), and set \(x_j=|q_j|/(\rho q_0)\).
The weighted Euler identity
\eqref{eq:uniform-euler} still holds. In particular
\(x_j\leq1\) for \(j\geq2\), and
\(\|x\|_2^2\leq1/(1-\rho^2)<2\).
The formal coefficient computed with these moduli is an
upper bound for the modulus of the actual coefficient by the
positive generalized binomial expansion.
This upper comparison uses only the finite jet of the formal
power.

The trial argument of Theorem~\ref{thm:explicit-small-p}
still implies \(x_1>c=3/4\). More precisely, now \(\rho\leq1/64\)
and
\(\rho/(1-\rho^2)\leq (257/256)\rho\).
The upper bound for the ratio to the same trial in the
case \(x_1\leq c\) is
\[
 h^n\exp\bigl((257/256)(e/h^2)n\rho\bigr)<1,
\]
using \(-\log h>1/16\) and
\(e/h^2=(16/9)e^{9/16}<32/9\).
The derivative comparison \eqref{eq:uniform-replacement} also
holds here: for \(j\geq2\), \(R_j/R\leq M_j\), and
\(\sum M_j\leq3n\rho\) follows from \(2\rho\leq1/32\).
Since \(R(x_1,0,\ldots,0)=1\), integrating these estimates
gives \(\log R(x)\leq3n\rho\,u\), where
\(u=(\sum_{j\geq2}x_j^2)^{1/2}\).
Relative to \(x_1\) alone, the log normalization decreases by
at least
\[
 \frac{s}{2}\log\left(1+
      \frac{\rho^2 u^2}{1+\rho^2 x_1^2}\right)
 \ \geq\ \frac{n u^2}{2(1+2\rho^2)}\ \geq\ \frac{nu^2}{3}.
\]
The \(x_1\)-only normalized objective attains its maximum at
\(x_1=(1-\rho^2)^{-1/2}\), and the value there is \(M(n,s)\).
Consequently
\[
 \log(C(n,p)/M(n,s))
 \leq 3n\rho\,u-nu^2/3\leq(27/4)n\rho^2.
\]
In particular the higher coordinates of any extremal factor
satisfy \(u\leq9\rho\).

Finally, the log ratio of \(M(n,s)\) to \((ns/e)^{n/2}/n!\)
is the sum of
\[
 \sum_{k=0}^{n-1}\log(1-k/s)
 \quad\hbox{and}\quad
 \frac{s}{2}\bigl((1-\rho^2)\log(1-\rho^2)+\rho^2\bigr).
\]
The first term lies between \(-n^2/s\) and \(0\).
The second lies between \(0\) and \(n^2/(2s)\), since its
bracket is the integral of \(-\log(1-v)\) from \(0\)
to \(\rho^2\) and \(-\log(1-v)\leq2v\) on this interval.
Combining these estimates proves the comparison.
\end{proof}

\section{Uniqueness on the inverse-degree scale}

\begin{theorem}\label{thm:linear-scale}
For every \(n\geq2\) and \(0<p\leq 1/(2048n)\) the normalized
maximizing factor is unique up to argument rotation and scalar
phase, and the value and factor are real analytic in the
interior of this range. On the convex set \(\Omega\) of
Theorem~\ref{thm:explicit-small-p} with the extra constraint
\(\|x_{\geq2}\|_2\leq9\rho\), the iteration of Corollary~\ref{cor:uniform-iteration} converges to this factor with contraction
factor at most \(11/12\).
\end{theorem}

\begin{proof}
Here \(s\geq4096n\) and \(\rho\le1/64\).
Theorem~\ref{thm:general-small-positive} lets us take the
coefficient moduli and optimize the formal power on the entire
sphere. Theorem~\ref{thm:uniform-comparison} puts every positive
normalized maximizer in the constrained convex set in the
statement and gives \(x_1>c=3/4\).
Use the same \(\Phi\) and its correction \(R\) as in
\eqref{eq:uniform-log-objective}.

We bound the Hessian of \(\psi=\log R\) more sharply than
in the previous section. For \(x_j>0\), regard each
multiplicity vector \(m=(m_2,\ldots,m_n)\) as a state, subject
to \(m_j\geq0\) and \(\sum_{j\ge2}j m_j\leq n\).
Set \(m_1=n-\sum_{j\ge2}j m_j\) and \(k=\sum_{j\ge1}m_j\).
Give each state probability equal to the corresponding
monomial weight in \(R\) divided by \(R\). Define
\[
 a_j(m)=
 \frac{(m_1)_{\underline{j}}}
      {x_1^j\rho^{j-1}\prod_{\ell=1}^{j-1}(s-k+\ell)}
 \qquad (j\geq2),
\]
where the falling factorial is zero when \(m_1<j\).
A birth in coordinate \(j\) has rate \(b_j(m)=x_j a_j(m)\).
A death has rate \(m_j\). The ratios of neighboring
monomial weights are \(b_j(m)/(m_j+1)\), so this finite
irreducible Markov chain is reversible for the chosen
probabilities.

We use the constants \(M_j\) from
\eqref{eq:uniform-replacement}. The numerator and denominator
estimates in that comparison give \(0\leq a_j\leq M_j\).
For a feasible adjacent pair \(m,m+e_i\), the numerator
decreases by at most \(ij n^{j-1}\), using telescoping
of the falling factorial. Each denominator factor rises
by at most \(i\), so the reciprocal product decreases
by at most a relative \(ij/(s-n)\). Hence
\[
 |a_j(m+e_i)-a_j(m)|\leq 2ij M_j/n .
\]
In the metric \(d(m,m')=\sum_{j\ge2}j|m_j-m'_j|\), the
Lipschitz constant of \(a_j\) is at most \(2j M_j/n\).
Indeed every pair is connected by a shortest path through
its coordinatewise minimum. Moreover birth rates are
nonincreasing along each increasing edge.

For adjacent states differing by \(e_i\), couple common
birth and death rates in the two copies. The one excess
death of coordinate \(i\) coalesces the states, and an
uncoupled birth in coordinate \(j\) increases distance by
at most \(j\). Thus the initial drift of their distance is
at most \(-\kappa i\), where
\[
 W=\sum_{j\ge2}j^2x_jM_j,\qquad \kappa=1-2W/n.
\]
The finite-state small-time bound extends by the triangle
inequality along shortest paths and iteration of the
semigroup, to give
\(\operatorname{Lip}(P_tf)\leq e^{-\kappa t}\operatorname{Lip}(f)\).
Explicitly, for a small time \(h\) the coupling gives a bound
\(1-\kappa h+O(h^2)\) for adjacent Lipschitz differences,
uniformly on the finite state space. Taking \(h=t/N\)
and iterating \(N\) times proves the claim.
By reversibility the variance identity for the chain is
\[
 \operatorname{Var}(f)=2\int_0^\infty
     \sum_j\mathbb E\!\left[
       b_j(m)\big(P_tf(m+e_j)-P_tf(m)\big)^2
                      \right]\,dt
 \ \leq\ \frac{W}{\kappa}\operatorname{Lip}(f)^2 .
\]
The identity follows by differentiating
\(\mathbb E[(P_tf)^2]\) and using convergence to its
stationary mean.

Coefficient replacement gives \(R_j/R=\mathbb E[a_j]\) and
\(R_{jk}/R=\mathbb E[a_j(m)a_k(m+e_j)]\); when \(m+e_j\)
is infeasible the term is zero. Consequently
\[
 \psi_{jk}=\operatorname{Cov}(a_j,a_k)
       +\mathbb E\!\left[a_j(m)(a_k(m+e_j)-a_k(m))\right].
\]
The variance and Lipschitz estimates give
\[
 |\psi_{jk}|
 \leq \left(2+\frac{4W}{\kappa n}\right)
          \frac{jk M_j M_k}{n}\qquad(j,k\geq2).
\]
On the constrained set \(x_j\leq9\rho\) for \(j\ge2\).
The geometric sums in the preceding proof therefore give
\(W\leq162n\rho^2\) and
\(S_2=\sum_{j\ge2}j(j+1)x_jM_j\leq162n\rho^2\).
Thus \(\kappa\geq1-324\rho^2>0\) and
\(4W/(\kappa n)<1\).

The identity \(x_1\psi_1=-\sum_{j\ge2}jx_j\psi_j\)
now bounds the remaining entries of the Hessian.
With \(T_1=\sum_{j\ge2}j M_j\leq6n\rho\), differentiating
this identity and summing absolute values gives
\[
 \|D^2\psi\|_{\rm op}
 \leq c^{-2}S_2+2c^{-1}T_1+\frac3n(T_1+c^{-1}W)^2
 \leq n\bigl(288\rho^2+16\rho+3(6\rho+216\rho^2)^2\bigr)
 <0.39n .
\]
For example the mixed terms use
\(x_1\psi_{1k}=-k\psi_k-\sum_{j\ge2}jx_j\psi_{jk}\).
The estimates extend to zero coordinates by continuity.
Adding the normalizing Hessian and \(n\log x_1\) gives
eigenvalues in \([-4n,-n/3]\) for \(D^2\Phi\), as in the
earlier contraction proof.
Thus there is at most one maximizing positive factor.
The jet identity makes it an interior stationary point
in the unconstrained real chart. The maximizing coordinates
are continuous in \(p\): every nearby maximizer also lies in
the fixed compact set \(\Omega\), and any limit maximizes the
continuous formal objective there at the base parameter.
Uniqueness identifies the only such limit. The negative
definite Hessian and the implicit function theorem now give
analytic dependence. The projection step is a contraction
with the indicated factor.
All complex Hardy maximizers lie in the same rotation
orbit by Theorem~\ref{thm:general-small-positive}.
\end{proof}

\section{The endpoint slope in every degree}

For a nonnegative trigonometric polynomial \(W\) with
\(\int_{\mathbb T}W\,dm=1\), put
\(\mathcal E(W)=\int_{\mathbb T}W\log W\,dm\), taking
\(0\log 0=0\). Let \(E_n\) be the maximum of \(\mathcal E(W)\)
over all such polynomials of degree at most \(n\). The following
result describes and computes the coefficient slope.

\begin{theorem}\label{thm:general-entropy}
For each \(n\geq1\),
\[
 \log C(n,p)=E_n(1/p-1)+O_n\bigl((1/p-1)^2\bigr)
 \qquad(p\uparrow1).
\]
Every limit of normalized extremal factors at \(p=1\) has a density
which achieves \(E_n\). Every density achieving \(E_n\) is the
squared modulus of a degree-\(n\) normalized polynomial with
\(n\) distinct unit-circle roots. Its polynomial factor is a
simple top eigenvector of the Toeplitz matrix with symbol \(\log W\).

The slope \(E_n\) has an exact finite computation by real
algebraic root isolation and logarithms. It is transcendental,
and \(E_{n+1}>E_n\). In particular,
\[
 2H_n-\frac{2n}{n+1}-\log(n+1)
 \ \leq\ E_n\ <\ \log(n+1),\qquad
 H_n=\sum_{j=1}^n\frac1j .
\]
For \(n\geq2\) the lower inequality is strict.
\end{theorem}

\begin{proof}
By Fej\'er--Riesz, the densities under consideration are precisely
\(W=|Q|^2\), where \(Q\) is a complex polynomial of degree at most
\(n\) and \(\|Q\|_2=1\). Compactness gives a maximizing density.
Let \(W_0\) be any such maximum. The function \(u\log u\) is
strictly convex on \([0,\infty)\). Its supporting-line inequality
at \(W_0(z)>0\), integrated on the circle, gives
\[
 \int_{\mathbb T}(W-W_0)\log W_0\,dm < 0
 \qquad\text{if }W\ne W_0.
\]
Here \(W_0\) has only finitely many zeros,
\(\log W_0\) is integrable, and equality of the two masses removes
the constant part of the derivative. Thus \(W_0\) is the unique
maximizing density for this linear functional.

On the space of polynomials of degree at most \(n\), let
\(T_n(\log W_0)\) be the Hermitian operator with quadratic form
\(\int |Q|^2 \log W_0\,dm\). Its largest eigenvalue is \(E_n\).
All its unit top eigenvectors have squared modulus \(W_0\).
The top eigenspace must be one-dimensional. Otherwise, choose
orthonormal vectors \(Q,R\) in this space and apply the squared
modulus condition to \((Q+R)/\sqrt 2\) and \((Q+iR)/\sqrt 2\).
It would follow that \(\overline Q R=0\) on the circle, which
is impossible for two nonzero polynomials.

Any normalized factor \(Q\) of \(W_0\) must therefore have degree
\(n\), since a lower-degree polynomial and its shift \(zQ\)
would give linearly independent top eigenvectors. Its roots must
lie on the unit circle. Reflecting a root \(\zeta\) off the circle
replaces \(z-\zeta\) by \(1-\bar\zeta z\), which preserves the
modulus on the circle and gives a different factor up to phase.
The same replacement at \(\zeta=0\) removes the factor \(z\).

The roots cannot repeat. If \(Q=(z-\zeta)^2R\), then on the
circle, outside the root,
\[
 (\zeta zR)\overline Q=-\left|\frac{Q}{z-\zeta}\right|^2 .
\]
Set \(M=T_n(\log W_0)-E_n I\). Since \(MQ=0\), testing it
against \(\zeta zR\) gives
\[
 \left\langle M\frac{Q}{z-\zeta},\frac{Q}{z-\zeta}\right\rangle=0 .
\]
This contradicts strict negativity of \(M\) off its kernel,
because \(Q/(z-\zeta)\) has degree \(n-1\).

Write \(\eta=1/p-1\), and let \(Q_p\) be an extremal factor.
Since \(W_p=|Q_p|^2\) is among the densities defining \(E_n\),
its entropy is at most \(E_n\). The \(H^1\) coefficient bound
\(C(n,p)\le\int W_p^{1+\eta}\,dm\) and the uniform expansion of
this integral give \(\log C(n,p)\le\eta E_n+O_n(\eta^2)\).
For the matching lower bound, take a normalized maximizing
endpoint factor \(Q_0\). Since all its roots lie on the circle,
it is self-reciprocal up to phase, hence
\(|[z^n]Q_0^2|=1\). The test function \(Q_0^2/\|Q_0^2\|_p\)
gives
\[
 \log C(n,p)\geq
 -p^{-1}\log\!\int |Q_0|^{2p}\,dm
 =\eta E_n+O_n(\eta^2).
\]
The Taylor estimates here are uniform over \(0\leq W\leq n+1\)
of mass one, since the second derivative of \(W^s\) is uniformly
bounded for \(s\) close to one. The same upper and lower estimates
show \(\mathcal E(|Q_p|^2)\to E_n\), and its limits thus attain
\(E_n\).

We give the algebraic computation explicitly. On the real
algebraic torus \(|\lambda_j|^2=1\) set
\[
 P(z)=\prod_{j=1}^n(1-\lambda_jz),\qquad U=\|P\|_2^2,
 \qquad s_k=\sum_{j=1}^n\lambda_j^k .
\]
Write \(P_k=[z^k]P\) and \(\mathbf P=(P_0,\ldots,P_n)^\mathsf T\).
Define \(T\) to have diagonal zero and entries \(T_{jk}=-s_{j-k}/(j-k)\)
for \(j>k\), with the other entries obtained by conjugation. The
Fourier expansion of \(\log |1-\lambda_j z|^2\) gives
\(T=T_n(\log |P|^2)\). Indeed, the expansion at radius \(r<1\)
is \(-\sum_{k\geq1}r^k(\lambda_j^k z^k+
\bar\lambda_j^k z^{-k})/k\), and its limit as \(r\uparrow1\)
holds in \(L^1\). For example, singularities near each zero can
be controlled by an integrable majorant proportional to
\(1+|\log|z-\bar\lambda_j||\) for \(r\geq1/2\).
The real algebraic system
\[
 |\lambda_j|^2=1\quad (1\leq j\leq n),\qquad
 T\mathbf P=A\mathbf P,\qquad A\in\mathbb R
\]
contains every endpoint optimizer. Each solution has
\[
 \mathcal E(|P|^2/U)=A-\log U .
\]
Moreover this value is constant along each connected smooth
cell of the solution set. To see this, differentiate it along
a piecewise continuously differentiable path in the cell, taking
\(Q=P/\sqrt U\). Differentiating the entropy under the integral
is justified by the bounded continuous first derivatives of
\(u\mapsto |u|^2\log|u|^2\). The equation
\(T_n(\log|Q|^2)Q=(A-\log U)Q\) and
\(d\|Q\|_2^2/dt=0\) therefore give zero derivative.

A real semialgebraic cell decomposition of the displayed
rational system provides a real algebraic sample on each cell.
The maximum of the finitely many values \(A-\log U\) so obtained
is exactly \(E_n\). Each can be compared with any other by
certified approximation, after an algebraic equality test.
Indeed, a tie would imply
\(A_1-A_2=\log(U_1/U_2)\), which by the Hermite--Lindemann
theorem forces \(A_1=A_2\) and \(U_1=U_2\). Since both
the constant and leading coefficients of \(P\) have modulus
one, \(U\geq2\). Thus every candidate value, and hence \(E_n\),
is transcendental.

Every degree-\(n\) maximizing density has degree exactly \(n\):
the top Fourier coefficient is a product of the nonzero constant
and leading coefficients of its factor. Thus \(E_{n+1}\) cannot
equal \(E_n\). The upper inequality follows from \(W\leq n+1\)
and is strict since equality would require \(W=n+1\) wherever
\(W>0\), incompatible with its mass. For the lower inequality,
take \(P=1+z+\cdots+z^n\). The coefficients through order \(n\)
of \(\log P\) are \(1/k\) in degree \(k\geq1\), and the
Fourier coefficients of the normalized squared modulus are
\(1-k/(n+1)\) for \(1\leq k\leq n\). Its entropy is therefore
\(2H_n-2n/(n+1)-\log(n+1)\).
For \(n\geq2\) this vector of constant coefficients is not
an eigenvector of its Toeplitz matrix \(T\), whose row sums
are \(H_j+H_{n-j}\) for \(0\leq j\leq n\), so this density
cannot achieve the maximum.
\end{proof}

\subsection{The entropy operator has spectral gap two}

\begin{proposition}\label{prop:entropy-gap}
For every entropy-maximizing density \(W=|Q|^2\) in degree
\(n\geq1\), the two largest eigenvalues of \(T_n(\log W)\)
are \(E_n\) and \(E_n-2\). Moreover all endpoint maximizers
with their factors normalized to have positive constant term
have the same constant term.
\end{proposition}

\begin{proof}
Choose the constant phase so that \(Q=Q^{*n}\).
The involution \(S:H\mapsto H^{*n}\) on the degree-\(n\)
coefficient space is conjugate-linear. The Hermitian Toeplitz
operator \(T=T_n(\log W)\) commutes with \(S\), as follows
directly from the entries depending on the difference of
their indices and Hermitian symmetry. For a real fixed vector
\(h=Sh\) orthogonal to \(Q\), the quantities
\(z^{-n/2}Q(z)\) and \(z^{-n/2}h(z)\) are real on the unit
circle, using either local square-root branch. Along
\(Q_t=(Q+th)/\sqrt{1+t^2\|h\|_2^2}\), the second derivative
of entropy at zero is consequently
\[
 2\langle (T-E_n I)h,h\rangle+4\|h\|_2^2 .
\]
For justification, the nearby self-reciprocal polynomials have
one zero near each simple unit root, and reciprocal symmetry
forces each such zero to stay on the unit circle.
The logarithms of their moduli then converge in \(L^1\) by
factorization and continuity of translations of
\(\log|1-e^{it}|\). The second derivative under the integral
along this real self-reciprocal path is therefore justified
by uniform integrability of these logarithms. Its extra term
is \(4\int(\Re(\bar Q h))^2/|Q|^2\,dm=4\|h\|_2^2\).

At a maximum this second derivative is nonpositive. The fixed
real subspace and its imaginary multiple span \(Q^\perp\).
Since \(T\) commutes with \(S\), the quadratic inequality on
the fixed subspace extends to \(Q^\perp\). This gives
\(\lambda_2(T)\le E_n-2\).
The rotation path
\(e^{-in\theta/2}Q(e^{i\theta}z)\) stays in the fixed subspace
and has constant entropy. Its tangent is
\(i(zQ'-nQ/2)\), a nonzero fixed vector orthogonal to \(Q\).
Orthogonality follows also from
\(|q_j|=|q_{n-j}|\). The constrained second derivative in
this direction is zero, giving equality in the quadratic
bound and hence \(\lambda_2(T)=E_n-2\).

For the constant term, use the algebraic endpoint system in
Theorem~\ref{thm:general-entropy}. On any maximizing cell its
algebraic sample points all have the same \(A,U\):
equality of \(A-\log U\) forces equality of both algebraic
quantities by Hermite--Lindemann. Algebraic points are dense
in each semialgebraic cell defined over the real algebraic
numbers, by real closed field transfer applied to rational
open boxes. Continuity makes the pair the same throughout the
maximizing cells. Thus every normalized positive constant
term is the same number \(U^{-1/2}\).
\end{proof}

\subsection{The motion of the zeros}

\begin{theorem}\label{thm:radial}
Let \(p\) tend to \(1\) from below along a sequence of normalized
extremal factors \(Q_p\) converging to \(Q_*\). Match each zero
\(\zeta_*\) of \(Q_*\) to the zero \(\zeta_p\) of \(Q_p\) that
converges to it. Then
\[
 \lim_{p\uparrow1}
 \frac{|\zeta_p|-1}{1/p-1}
 =
 \frac{\|Q_*/(z-\zeta_*)\|_2^2}{|Q_*'(\zeta_*)|^2}.
\]
In particular, for each fixed degree, the distances of all
extremal zeros from the circle are bounded above and below by
positive constants times \(1/p-1\) when \(p\) is sufficiently
close to \(1\). Every limiting velocity belongs to
\[
 \left[\,\frac1n,\ \log(n+1)+2n\log2\,\right].
\]
\end{theorem}

\begin{proof}
Rotate and normalize the factors so that \([z^n]Q_p^s=C(n,p)>0\),
passing to a subsequence of rotations if necessary. The quotient
in the theorem is unchanged by rotation and multiplication by
unit constants. Put \(\tau=1/p-1\), \(C=C(n,p)\), and use the
polynomials \(R,E\) from the reflected boundary identity. At a
zero \(\zeta_p\) of \(Q_p\), both \(R\) and \(E^{*2n}\) vanish,
the latter by its factor \(Q_p\). The polynomial version of
\eqref{eq:boundary-identity} therefore gives
\[
 (1+\tau)\,Q_p^{*n}(\zeta_p)\zeta_p Q_p'(\zeta_p)
   = \tau\,\zeta_p^n D_{|Q_p|^2}(\zeta_p).
\]
On the right \(D_W\) denotes the Laurent polynomial obtained
from \(W\) by the multiplier \(|k|\).
By Theorem~\ref{thm:general-entropy}, the limiting factor has
simple unit roots. Thus there is exactly one zero near
\(\zeta_*\) for small \(\tau\), and \(Q_p^{*n}\) has a zero at
\(\zeta_p^\sharp=1/\overline{\zeta_p}\). Since
\((\zeta_p-\zeta_p^\sharp)/(|\zeta_p|-1)\to2\zeta_*\),
\[
 \frac{Q_p^{*n}(\zeta_p)}
      {|\zeta_p|-1}
 \longrightarrow 2\zeta_*(Q_*^{*n})'(\zeta_*)
 =-2\zeta_*^{n-1}\overline{Q_*'(\zeta_*)}.
\]
The last identity is coefficient reversal differentiated at
a unit-circle zero. Also
\(D_{|Q_*|^2}(\zeta_*)=-2\|Q_*/(z-\zeta_*)\|_2^2\)
by the same \(|k|\)-multiplier calculation as in the boundary
proof. Divide the evaluated identity by the convergent factors
and cancel \(-2\zeta_*^n\) to obtain the claimed velocity.

The set of limiting optimizer pairs \((Q_*,\zeta_*)\) is
contained in the compact set of normalized entropy-maximizing
factors and their roots. On that compact set the displayed
velocity is continuous and strictly positive by simplicity.
A subsequence argument gives the asserted uniform bounds.
For the numerical lower bound, apply Cauchy--Schwarz to
the evaluation of \(Q_*/(z-\zeta_*)\) at \(\zeta_*\),
a polynomial of degree \(n-1\). For the upper bound,
sum the velocities over all roots. Since \(d=[z^n](Q_p/Q_p(0))\),
\(-\log|d|=\sum_j\log|\zeta_{p,j}|\). The stationary value
and the endpoint asymptotic give
\[
 \sum_j\frac{\|Q_*/(z-\zeta_{*,j})\|_2^2}
                  {|Q_*'(\zeta_{*,j})|^2}
     =E_n-2\log|Q_*(0)|.
\]
This is at most \(\log(n+1)+2n\log2\) by the reciprocal-root
bound on the constant term.
\end{proof}

\subsection{A global consequence of the endpoint slope}

\begin{proposition}
For all \(n\geq1\) and \(0<p<1\),
\[
 C(n,p)<\exp\bigl(E_n(1/p-1)\bigr)
 \leq \left(\csc\frac{\pi}{2(n+1)}\right)^{1/p-1}.
\]
The constant \(E_n\) in the first exponential bound is optimal.
\end{proposition}

\begin{proof}
At a Hardy extremizer, \(|Q|^2\) is strictly positive on
\(\T\), whereas all entropy-maximizing factors have unit roots.
Thus every slope from Proposition~\ref{prop:slopes} is strictly
less than \(E_n\). The finite analytic partition in
Theorem~\ref{thm:branches} and continuity at \(\tau=0\) give
\(J_n(\tau)<\tau E_n\) by integration. Sharpness of \(E_n\)
follows from its endpoint limit.

For the explicit bound put \(m_0=n+1\), and sample any
\(H^2\)-unit polynomial \(Q\) at two interlaced equispaced
grids. In the first grid let
\[
 u_\theta(j)=m_0^{-1/2}
      Q\!\left(e^{i(\theta+2\pi j/m_0)}\right),
 \qquad 0\leq j<m_0 .
\]
This is a unit vector by discrete Fourier orthogonality.
The map from \(u_\theta\) to \(u_{\theta+\pi/m_0}\) is unitary
and its entries have magnitude at most
\(c=(m_0\sin(\pi/(2m_0)))^{-1}\) by a geometric sum.
For a unitary matrix \(V\) with entries so bounded, interpolation
between its \(\ell^1\)-to-\(\ell^\infty\) norm and its
\(\ell^2\)-to-\(\ell^2\) norm gives
\(\|Vu\|_{r'}\leq c^{2/r-1}\|u\|_r\) for \(1\leq r\leq2\)
and \(1/r+1/r'=1\). Differentiating at \(r=2\) from the
left gives the Shannon entropy inequality
\[
 H(|u|^2)+H(|Vu|^2)\geq-2\log c,\qquad
 H(v)=-\sum_j v_j\log v_j .
\]
Indeed the derivative of \(\log\|u\|_r\) at 2 is
\(-H(|u|^2)/4\), and \((r')'(2)=-1\).
Averaging over \(\theta\),
\[
 \int |Q|^2\log |Q|^2\,dm
   =\log m_0-\operatorname{avg}_{\theta}H(|u_\theta|^2)
   \leq-\log\sin(\pi/(2m_0)).
\]
Taking the maximum over unit polynomials proves the result.
\end{proof}

\subsection{Cancellation of the quadratic term}

In fact, the error in Theorem~\ref{thm:general-entropy} improves to
\begin{equation}\label{eq:general-hardy-endpoint-cubic}
 \log C(n,p)=E_n(1/p-1)+O_n\bigl((1/p-1)^3\bigr).
\end{equation}
The cancellation does not require uniqueness of an entropy-maximizing
density.

Put \(\delta=2/p-2\) and let \(J\) be conjugate reversal in degree
\(n\). Multiplication by a phase makes \([z^n]Q^2>0\) for all
normalized factors under consideration near \(p=1\).
In this gauge write
\[
 Q=(1-\|y\|_2^2)^{1/2}X+i y,\qquad
 JX=X,\quad Jy=y,\quad \|X\|_2=1,\quad \langle X,y\rangle=0.
\]
Indeed the real and imaginary parts for the real form fixed by
\(J\) have real inner products, and the phase condition on
\([z^n]Q^2\) makes them orthogonal. For an extremal factor, the stationary value and the upper
bound \(C(n,p)\leq(n+1)^{\delta/2}\) imply
\(-\log|[z^n](Q/Q(0))|=O_n(\delta)\); here
\(|Q(0)|\geq2^{-n}\). The leading coefficient modulus is the
product of the moduli of the reciprocal roots. Thus their
radial distances from the circle sum to \(O_n(\delta)\).
Projecting the reciprocal roots radially to the circle and
telescoping the factored polynomial shows that its normalized
factor differs by \(O_n(\delta)\) in coefficient norm. In this
gauge this gives
\(\|y\|_2=O_n(\delta)\). To see this after the phase choice,
the projected factor is within \(O_n(\delta)\) beforehand, its
coefficient of the squared factor has modulus one, and the
phase correction to make this coefficient positive preserves
the estimate. The normalized projected factor is then fixed by
\(J\). The resulting \(X\)'s approach the compact set of
entropy maximizers in this real form. They have unit-circle
roots for \(p\) close to one. Indeed each of these limiting
polynomials is simple, and nearby self-reciprocal polynomials
have each of their roots on the circle by local uniqueness of
the root and reciprocal symmetry.

For a self-reciprocal \(X\) in this neighborhood take an analytic
logarithm \(L=\log X\) near zero, and put
\[
 F=\Pi_n(XL)=u+i v,\qquad Ju=u,\quad Jv=v ,
\]
where \(\Pi_n\) is coefficient projection onto degrees \(0,\ldots,n\).
The subscript \(\perp\) denotes the orthogonal component to \(X\).
With the inner product linear in the first factor, define
\(\mu=\langle F,X\rangle\). Formal coefficient multiplication and
\(JX=X\) give
\[
 [z^n]X^2L=\mu,\qquad
 [z^n]X^2L^2=[z^n]F^2=\langle F,JF\rangle.
\]
The first equality truncates only terms of degree above \(n\).
Also \(\Re\mu=\mathcal E(|X|^2)/2\).
In particular,
\[
 \Re\bigl([z^n]X^2L^2-\mu^2\bigr)
       =\|u_\perp\|_2^2-\|v_\perp\|_2^2 .
\]

The function \(\Phi(Q,\delta)=\log|[z^n]Q^{2+\delta}|\)
is real analytic in the coefficients and in \(\delta\) locally
where the two relevant coefficients \(Q(0)\) and \([z^n]Q^{2+\delta}\)
are nonzero. This follows from the finite binomial formula
through order \(n\). The Taylor expansion in the displayed
normal coordinates, uniformly when \(\|y\|=O_n(\delta)\), is
\begin{align*}
 \Phi(Q,\delta)
 &=\frac{\delta}{2}\mathcal E(|X|^2)
   -2\|y\|_2^2-2\delta\langle v_\perp,y\rangle\\
 &\quad+\frac{\delta^2}{2}
          \bigl(\|u_\perp\|_2^2-\|v_\perp\|_2^2\bigr)
          +O_n(\delta^3)\\
 &=\frac{\delta}{2}\mathcal E(|X|^2)
   +\frac{\delta^2}{2}\|u_\perp\|_2^2
   -2\bigl\|y+\delta v_\perp/2\bigr\|_2^2
   +O_n(\delta^3).
\end{align*}
Here \(\Phi(Q,0)=\log(1-2\|y\|_2^2)\).
The mixed term comes from differentiating \(Q^2\log Q\)
in the direction \(i y\), which gives
\(i Xy(2L+1)\). Its real coefficient at degree \(n\) is
\(-2\langle v_\perp,y\rangle\), since \([z^n]Xy=0\).
The pure \(\delta^2\) term is the identity for \(\mu\) and \(F\)
above. Uniformity follows from compactness and a finite cover
by logarithm branches.

On the unit sphere in the real form fixed by \(J\), the entropy
near the maximizing profiles is real analytic, by the formula
\(\mathcal E(|X|^2)=2\Re[z^n]X^2\log X\).
Its spherical gradient is \(4u_\perp\), since
\(T_n(\log|X|^2)X=F+JF\).
There is a constant \(K_n\) on a neighborhood of the compact
maximizing set such that
\(\|u_\perp\|_2^2\leq K_n(E_n-\mathcal E(|X|^2))\).
For completeness, step from \(X\) in the unit gradient direction
on the sphere. A uniform Hessian bound gives an increase at least
\(\|\nabla\mathcal E\|_2^2/(2M)\) at step
\(\|\nabla\mathcal E\|_2/M\), with \(M\) enlarged so this step
stays in the neighborhood. The increase is at most the
difference from the global maximum. Thus for sufficiently small
\(\delta>0\) the expansion gives the required upper bound
\(\Phi(Q,\delta)\leq\delta E_n/2+O_n(\delta^3)\).

For the lower bound, fix a sequence of extremal factors tending
to an entropy maximizer \(X_*\) in the chosen phase gauge.
The first-order jet along this sequence, expanded through
degree \(n\), gives
\[
 X+i y+\delta F=(1+\delta\mu)(X-i y)+O_n(\delta^2).
\]
Here \(X,y,F,\mu\) refer to the sequence at exponent \(2+\delta\).
It follows that \(y/\delta\) tends to \(-v_{*,\perp}/2\), since
\(F_*-\mu_* X_*=i v_{*,\perp}\) at the entropy maximum.
Let \(\widehat\zeta\) be a zero of \(X\) near a simple zero
\(\zeta_*\) of \(X_*\). Then \(|\widehat\zeta|=1\). The root of
\(Q\) near it satisfies
\[
 \zeta-\widehat\zeta
 =-\frac{i y(\widehat\zeta)}{X'(\widehat\zeta)}+O_n(\delta^2).
\]
By Theorem~\ref{thm:radial}, the limiting outward radial speed
per unit \(\delta\) of this displacement is strictly positive.
Therefore the fixed-profile trial
\(Q_{\mathrm{trial}}=\sqrt{1-\|y_*\|_2^2}\,X_*+i y_*\), with
\(y_*=-\delta v_{*,\perp}/2\), has strictly exterior roots for
all sufficiently small positive \(\delta\): its first-order
radial speed equals that of the sequence, by the preceding root
differentiation formula. Its power is thus admissible. In its
Taylor expansion \(u_{*,\perp}=0\) and the square term is zero,
so its value proves the matching lower bound in
\eqref{eq:general-hardy-endpoint-cubic}.

\section{Scope and literature}

The determinant in Theorem~\ref{thm:matrix-all-n} uses an
explicit finite multiplication matrix for the original Hardy
coefficient norm at every exponent. Its size is a uniform
upper bound and can be large. In degree five, the distinguished
root of a degree-fifteen eliminant gives a much smaller formula
and unique extremizers for the full exponent range \cite{Fifth}.
The arbitrary-degree determinant, finite branch theorem,
and endpoint conclusions do not assert uniqueness for every
exponent in every degree. The earlier cell-based calculation
\cite{Low} already provided an exact representation. The
determinant replaces that sampling step and gives a uniform
degree bound. The entropy maximum here concerns a different
extremal problem from the entropy minimum in the
Agler--McCarthy program \cite{Agler,LeiZhang}.

\small
\begin{thebibliography}{99}
\setlength{\itemsep}{2pt}
\bibitem{BCR}
J.~Bochnak, M.~Coste and M.-F.~Roy, \emph{Real Algebraic Geometry},
Springer, 1998.

\bibitem{BrevigSaksman}
O.~F.~Brevig and E.~Saksman, Coefficient estimates for \(H^p\)
spaces with \(0<p<1\), \emph{Proc. Amer. Math. Soc.} \textbf{148}
(2020), 3911--3924.
\href{https://doi.org/10.1090/proc/14995}{DOI: 10.1090/proc/14995}.

\bibitem{BrevigSovik}
O.~F.~Brevig and \O.~S{\o}vik, Extremals for linear functionals on
\(H^p\) are outer when \(0<p<1\),
\href{https://arxiv.org/abs/2610.21788}{arXiv:2610.21788}, 2026.

\bibitem{Low}
\emph{Sharp third and fourth Taylor coefficients in Hardy spaces
below \(H^1\)}, preceding manuscript, November 2026.

\bibitem{High}
\emph{Sharp fifth and sixth Taylor coefficients in Hardy spaces
below \(H^1\)}, preceding manuscript, November 2026.

\bibitem{Fifth}
\emph{The fifth Taylor coefficient norm in Hardy spaces below
\(H^1\)}, preceding manuscript, November 2026.

\bibitem{Agler}
J.~Agler and J.~E.~McCarthy, The Krzy\.z conjecture and an
entropy conjecture, \href{https://arxiv.org/abs/1803.09718}{arXiv:1803.09718}.

\bibitem{LeiZhang}
J.~Lei and T.~Zhang, Proof of the Agler--McCarthy entropy conjecture,
\href{https://arxiv.org/abs/2605.03949}{arXiv:2605.03949}, 2026.

\bibitem{VMRC}
Verifiable Mathematical Reasoning Corpus, manuscripts [248] and
[249] on Hardy coefficient problems,
\href{https://github.com/QEEDDD/Verifiable-Mathematical-Reasoning-Corpus/tree/e73bae9fd0c319129b19aacc5c317f7abf2865ef}{snapshot e73bae9f},
2026.
\end{thebibliography}
\end{document}
